\chapter{RESULTADOS}

\section{Configuración experimental}

La validación del modelo propuesto se realizó sobbre el conjunto de datos CASIA-B. A diferencia de aproximaciones tradicionales basadas en imágenes individuales, el presente trabajo adopta un enfoque espacio-temporal, en donde la marcha es modelada bajo el timepo. \\
En este contexto, cada muestra de entrada no corresponde a una imagen, sino a un volumen de datos compuestos por $T=24$ fotogramas consecutivos. Este volumen se representa como un tensor de dimensión $[T, C, H, W]$. \\
Este paradigma permite a la red aprender patrones cinemáticos complejos del ciclo de marcha, en lugar de depender únicamente de información estática de apariencia. 


\subsection{Protocolo de evaluación}

Los experimentos fueron diseñados bajo las siguientes condiciones:

\begin{itemize}
    \item \textbf{Dataset:} Se utilizaron los 124 sujetos oficiales de CASIA-B.
    
    \item \textbf{Condiciones de evaluación:} Se consideraron las tres condiciones estándar del dataset:
    \begin{itemize}
        \item Caminata normal (\texttt{NM})
        \item Caminata con bolsa (\texttt{BG})
        \item Caminata con abrigo (\texttt{CL})
    \end{itemize}
    
    \item \textbf{Entrada del modelo:} Secuencias de $T=24$ fotogramas consecutivos, formando tensores volumétricos $[T, C, H, W]$.
    
    \item \textbf{Pre-procesamiento:} Siluetas normalizadas en el rango [0,1] y redimensionadas a $64\times64$ píxeles.
    
    \item \textbf{Métricas de evaluación:} Comparaciones bajo \textbf{Rank-1 Accuracy} y \textbf{Mean Average Precision (mAP)}.
\end{itemize}

\section{Resultados cuantitativos}

El modelo espacio-temporal propuesto obtuvo los siguientes resultados sobre CASIA-B:

\begin{itemize}
    \item \textbf{mAP global:} 49.55\%
    \item \textbf{Rank-1 (Normal - NM):} 99.5\%
    \item \textbf{Rank-1 (Bolsa - BG):} 76.0\%
    \item \textbf{Rank-1 (Abrigo - CL):} 49.6\%
\end{itemize}

Estos resultados evidencia una mejora significativa respecto a otros enfoques basados en representaciones estáticas, especialmente en términos de la métrica mAP, lo cual indica una mejor estructura del espacio de \textit{embeddings}.

\section{Interpretación del desempeño}

El alto desempeño obtenido en la condición de caminata normal (\textbf{99.5\% Rank-1}) nos confirma que el modelo tiene la capacidad de capturar con precisión la dinámica de la marcha cuando no hay perturbaciones externas. Este resultado es consistente con la representación volumétrica permite modelar directamente la cinemática del movimiento humano. \\
Por otro lado, el desempeño en condiciones más desafiantes revela limitaciones importantes:

\begin{itemize}
    \item En la condición \textbf{BG (76.0\%)}, la presencia de objetos externos introduce ruido en la silueta, afectando al patrón temporal.
    
    \item En la condición \textbf{CL (49.6\%)}, el abrigo afecta significativamente la geometría del cuerpo, ocultando información clave de la marcha, lo cual impacta negativamente en la discriminabilidad del modelo.
\end{itemize}

Por otro lado, el valor de \textbf{mAP (49.55\%)} representa una mejora frente a configuraciones previas, indicando que el modelo no solo acierta en la primera posición, sino que mantiene su valor en un ranking más amplio.

\section{Análisis del Modelado Temporal}

La importancia del modelo radica en su capacidad para transformar la información temporal en una representación geométrica discriminativa. A diferencia de enfoques basados en imágenes individuales, el modelo procesa secuencias completas como volúmenes tridimensionales, donde la dimensión temporal es tratada como una variable explícita del sistema. \\
El módulo de \textit{temporal pooling} desempeña un rol clave en este proceso, al condensar la información temporal mediante operaciones de promedio (GAP) y máximo (GMP), lo que permite capturar el comportamiento global y otros eeventos más discriminativos del ciclo de marcha. \\
En este sentido, el modelo no aprende solo de la forma de la silueta, sino también sobre la evolución temporal de la misma.

\section{Comparación con el Estado del Arte}

\begin{table}[h]
\centering
\caption{Comparación de Rank-1 (\%) por condición}
\begin{tabular}{lccc}
\hline
\textbf{Método} & \textbf{NM} & \textbf{BG} & \textbf{CL} \\
\hline
GaitGL & $\sim$98 & $\sim$95 & $\sim$87 \\
GaitTAKE & 98.0 & 97.5 & 92.2 \\
\hline
\textbf{Modelo Propuesto} & \textbf{99.5} & 76.0 & 49.6 \\
\hline
\end{tabular}
\end{table}

El modelo propuesto supera a métodos bajo la condición de caminata normal. Sin embargo, también presenta una degradación considerable en condiciones de variación de apariencia. Esta diferencia se explica principalmente por el tipo de información utilizada:

\begin{itemize}
    \item Los métodos del estado del arte integran información adicional como pose o mecanismos de atención temporal.
    
    \item El modelo propuesto utiliza únicamente siluetas, lo cual limita su capacidad para manejar cambios estructurales severos como el uso de abrigo.
\end{itemize}

No obstante, es importante destacar que el modelo logra resultados competitivos en condiciones controladas sin la implementación de arquitecturas adicionales, resaltando su eficiencia.

\section{Discusión}

Los resultados obtenidos evidencian que el modelado espacio-temporal propuesto basado en siluetas es efectivo para capturar patrones de marcha en condiciones ideales. Sin embargo, la falta de información estructural limita su robustez frente a variaciones de apariencia. \\
Este comportamiento es consistente con la literatura, donde se ha demostrado que la incorporación de información de pose u otros mecanismos de atención mejora significativamente el desempeño bajo condiciones complejas. \\
En este sentido, la presente investigación establece una base para futuras extensiones, como la integración de características de pose o el uso de mecanismos de atención temporal, con el objetivo de mejorar la consistencia del modelo ante perturbaciones.

\section{Conclusiones del Capítulo}

En resumen, los resultados demuestran que:

\begin{itemize}
    \item El modelado volumétrico permite capturar eficazmente la dinámica de la marcha.
    \item El modelo alcanza desempeño sobresaliente en condiciones normales.
    \item Persisten desafíos importantes en escenarios con variaciones de apariencia.
\end{itemize}

Los resultados validan que la incorporación explícita de la dimensión temporal es un factor clave para mejorar el reconocimiento por marcha.